\begin{afterword}
\summaryhead

The post answers readers troubled by many-worlds: by the thought of constantly splitting into other people, and by whether buckling a seat belt here makes another self unbuckle. Its answer is ``Egan's Law'': it all adds up to normality. Every new physical theory must reproduce the successes of the old one, and many-worlds is not there to make strange predictions; the reader has always lived in this universe.

Decisions have no special link to branching, which comes from decoherence and happens all the time, even to rocks. Choices in one world do not affect another. Worlds outside your future can be treated as a ``generalized past'', beyond your responsibility. A 90\% chance and 9 worlds in 10 call for the same decision, and very improbable worlds deserve attention only in proportion to their weight; a die rolled until it shows 9 twelve times running gives a feel for one in a trillion. The post allows two ethical exceptions, that average utilitarianism looks better and that being ``first'' to know something means nothing, and otherwise advises adopting no strange belief because of many-worlds.

\responsehead

Most of the practical advice is sound, and it does not depend on the interpretation. Decisions do not cause branching; choices in one world do not reach another; very unlikely outcomes deserve attention only in proportion to their likelihood; a shaky grasp of physics is no reason for a strange belief. On single-world readings of quantum mechanics the same advice holds, and the worries it answers do not arise. The die makes one in a trillion concrete, and its arithmetic checks.

Two claims are stated more strongly than the evidence allows. Splitting is called the ``straightforward and unavoidable prediction of quantum mechanics''. It is what the many-worlds reading says the theory describes; the single-world readings predict the same results in every experiment done so far, and the book's argument for many-worlds comes later, in ``Many Worlds, One Best Guess''. The claim that seeing another world is ``unambiguously prohibited outright by the laws'' is stronger than the author's own earlier statement that the interaction between branches ``may never actually become zero''. On the standard account decoherence suppresses such interference rather than forbidding it; the practical point survives.

One step assumes what the author elsewhere calls an open problem. That a 90\% chance and ``9 out of 10 worlds'' call for the same decision requires branch weights to act as probabilities. Five weeks earlier, in ``The Born Probabilities'', the author wrote of that question: ``I don't know. It's an open problem.'' The attempt to settle it by Deutsch and Wallace remains contested.

The first ethical exception rests on a premise the post does not state. The case for average utilitarianism, that many people already exist in other worlds, persuades only someone who already thinks an added person is worth less when many exist; on the total view the other worlds change nothing. The second, about being ``first'' to know, repeats ``Joy in Discovery'' and needs only a large universe, not many worlds.

In short: sound practical advice that holds on any reading of quantum mechanics, with the many-worlds reading stated as settled, one step that assumes an answer to what the author called an open problem, and a case for average utilitarianism that rests on an unstated premise.
\end{afterword}
